\chapter{Problema de Pesquisa}
O interesse pela robótica autônoma e sistemas inteligentes surgiu ao longo da minha formação acadêmica, consolidando-se especificamente durante o desenvolvimento do projeto "Atom". Minha trajetória de iniciação científica foi marcada pela construção experimental de um manipulador robótico antropomórfico, onde participei ativamente de todas as etapas, desde a manufatura aditiva (impressão 3D) baseada na plataforma InMoov até a montagem eletromecânica e integração inicial de hardware. Durante esse processo, documentado no projeto de Grupo de Iniciação Científica (GIC), deparei-me com os desafios práticos de integrar a lógica de controle com a mecânica de precisão. No entanto, a experiência revelou uma lacuna significativa: embora tivéssemos um "corpo" funcional, o robô carecia de "visão" e inteligência para operar sem intervenção humana direta. Essa inquietação motivou o aprofundamento em estudos sobre desempenho de linguagens de programação para visão computacional, conforme explorado em pesquisas anteriores sobre o "Custo da Abstração". O presente projeto é, portanto, a evolução natural dessa trajetória, buscando unir a base mecânica já construída à inteligência artificial embarcada.

O manipulador robótico "Atom", desenvolvido anteriormente, possui estrutura mecânica funcional baseada no projeto InMoov, mas opera sem autonomia, dependendo de controle manual ou sequências pré-gravadas. O problema central é a limitação de processamento na borda: Redes neurais tradicionais de detecção de objetos exigem alto poder computacional (GPUs).

No contexto socioeconômico contemporâneo, a difusão da Indústria 4.0 e da automação inteligente tem ampliado a demanda por soluções autônomas de baixo custo, especialmente relevantes para países em desenvolvimento e setores com restrições orçamentárias (Schwab, (2016) \cite{Schwab2016} ).